\chapter{Requirements and Specification}

This chapter defines the functional and non-functional requirements for the synthetic ICU trajectory generation system, derived from the research objectives stated in Chapter~1. Requirements are prioritised using the MoSCoW method (Must, Should, Could, Won't) and linked to verification criteria used in later evaluation.

%==============================================================================
\section{Use Cases}
%==============================================================================

Synthetic ICU trajectory generation is a research-facing capability rather than a consumer product. Accordingly, the actors identified here are organisations and individuals engaged in clinical data science, machine-learning research, and healthcare data governance---the three communities for which synthetic EHR generators have the most documented practical value \cite{yan2022benchmarking,kaabachi2025scoping}.

\subsection{Actors}

\begin{description}
    \item[\textbf{Clinical Data Custodian (CDC).}] A data engineer or clinical informatics specialist at a hospital or research institute who holds a credentialed dataset (e.g.\ MIMIC-III) and wishes to derive a shareable synthetic analogue. This actor is primarily motivated by data-governance obligations---privacy regulations such as HIPAA and GDPR restrict direct sharing of raw patient records \cite{dataProtectionAct2018}. The CDC runs the training and generation pipeline and is responsible for evaluating the privacy properties of the synthetic output before any release.

    \item[\textbf{Downstream ML Researcher (MLR).}] A machine-learning researcher or data scientist who requires realistic ICU time-series data to develop, benchmark, or fine-tune clinical prediction models (e.g.\ mortality prediction, sepsis onset, ICU length-of-stay classifiers) but does not have direct access to a credentialed patient dataset. The MLR receives a synthetic cohort from the CDC and uses it as a training corpus \cite{yan2022benchmarking}.

    \item[\textbf{Privacy Auditor (PA).}] A member of an institutional review board, data governance committee, or responsible-disclosure team who must verify that a synthetic dataset does not memorise training-set individuals before it is approved for external release. The PA runs the privacy evaluation suite and interprets the resulting risk metrics against institutional thresholds \cite{kaabachi2025scoping,yao2025dcrdelusion}.
\end{description}

\subsection{Use Case 1: Privacy-Preserving Data Release}

\textbf{Actor:} Clinical Data Custodian (CDC). \\
\textbf{Goal:} Generate a synthetic ICU cohort that can be shared with external collaborators without disclosing real patient records.

\begin{enumerate}
    \item The CDC obtains credentialed MIMIC-III access and produces hourly ICU tables via a MIMIC-Extract style preprocessing pipeline.
    \item The CDC runs the data preparation scripts (\texttt{00\_prepare}, \texttt{01\_preprocess}) to build autoregressive training tables.
    \item The CDC fits the tabular preprocessor on the full dataset (\texttt{01b\_fit\_preprocessor}) to capture all categorical values and numeric ranges.
    \item The CDC trains the generative model (\texttt{02\_train\_model}) using the preferred backbone (ForestFlow or HS3F).
    \item The CDC generates a synthetic patient cohort (\texttt{03\_generate\_synthetic}) at the desired scale.
    \item The CDC runs the privacy evaluation suite and reviews DCR diagnostics and membership inference attack scores.
    \item If privacy risk is within acceptable thresholds, the synthetic cohort is packaged and shared with the MLR. Real patient tables are never distributed.
\end{enumerate}

\textbf{Success criterion:} DCR overfitting-protection score $\geq$ threshold; DOMIAS-style MIA ROC-AUC near chance level; no exact-match records in synthetic output.

\subsection{Use Case 2: Synthetic-Data-Augmented Model Development}

\textbf{Actor:} Downstream ML Researcher (MLR). \\
\textbf{Goal:} Train and evaluate clinical prediction models using a synthetic cohort as a substitute or augmentation for real ICU data.

\begin{enumerate}
    \item The MLR receives a synthetic ICU trajectory CSV from the CDC (Use Case~1).
    \item The MLR trains a downstream clinical classifier (e.g.\ mortality predictor, LOS classifier) on the synthetic training split.
    \item The MLR evaluates the trained classifier on a held-out \emph{real} test set (TSTR protocol \cite{esteban2018rcgan}).
    \item The MLR compares TSTR performance to a real-data baseline (TRTR) to assess the utility of the synthetic cohort.
    \item If utility is insufficient, the MLR requests a regeneration from the CDC with adjusted hyperparameters or a larger synthetic cohort.
\end{enumerate}

\textbf{Success criterion:} Synthetic-trained classifier achieves ROC-AUC within an acceptable gap of the real-trained baseline; F1 scores indicate minority-class retention.

\subsection{Use Case 3: Privacy Audit Prior to Data Release}

\textbf{Actor:} Privacy Auditor (PA). \\
\textbf{Goal:} Independently verify that the synthetic dataset approved by the CDC does not expose training-set individuals before external release is approved.

\begin{enumerate}
    \item The PA receives the synthetic cohort and the evaluation artifacts (sweep log, privacy metrics) from the CDC.
    \item The PA re-runs the privacy evaluation module (\texttt{privacy\_metrics.py}) on the synthetic cohort against the held-out real data.
    \item The PA inspects DCR distributions, baseline and overfitting protection scores, and MIA ROC-AUC.
    \item The PA cross-references findings against the DCR-delusion literature \cite{yao2025dcrdelusion}---recognising that proximity metrics alone are not sufficient anonymity evidence---and supplements with DOMIAS-style attack results.
    \item The PA issues a release recommendation, with caveats on subgroup populations or features where risk remains elevated.
\end{enumerate}

\textbf{Success criterion:} MIA attacker advantage is near zero; no individual synthetic record matches a real training record exactly; privacy report is documented alongside the synthetic dataset.

%==============================================================================
\section{Functional Requirements}
%==============================================================================

\begin{table}[htbp]
\centering
\small
\setlength{\tabcolsep}{3pt}
\begin{tabularx}{\textwidth}{@{}ll>{\raggedright\arraybackslash}Xl@{}}
\toprule
\textbf{ID} & \textbf{Priority} & \textbf{Requirement} & \textbf{Verification} \\
\midrule
FR1 & Must & The system must preprocess MIMIC-III hourly ICU tables into an autoregressive format with lag-1 features, separating static and dynamic columns. & Unit tests; output CSV inspection \\
\midrule
FR2 & Must & The system must fit a tabular preprocessor on the full dataset via streaming to capture all categorical values and numeric ranges. & Preprocessor artifact; inverse-transform round-trip \\
\midrule
FR3 & Must & The system must train a conditional generative model (ForestFlow or HS3F) on the autoregressive data, learning $p(\mathbf{x}_t \mid \mathbf{x}_{t-1}, \mathbf{s})$. & Model artifact; unit tests \\
\midrule
FR4 & Must & The system must generate synthetic patient trajectories by autoregressively sampling from the trained model. & Output CSV; row count \\
\midrule
FR5 & Must & The system must evaluate synthetic data quality via KS statistics, correlation preservation, and clinical range violation rates. & Quality metrics output \\
\midrule
FR6 & Must & The system must evaluate downstream utility via multi-task patient-level TSTR covering mortality and ICU length-of-stay prediction. & TSTR metrics in sweep CSV \\
\midrule
FR7 & Should & The system should train an IID hour-0 model for fully synthetic generation without requiring real initial conditions. & Hour-0 model artifact \\
\midrule
FR8 & Should & The system should evaluate privacy risk via DCR diagnostics and membership inference attacks. & Privacy metrics output \\
\midrule
FR9 & Should & The system should support both ForestFlow and HS3F backbones via a unified interface with backbone selection. & CLI flag; unit tests \\
\midrule
FR10 & Could & The system could provide visualisation of sweep results across hyperparameter configurations. & Plot output files \\
\bottomrule
\end{tabularx}
\caption{Functional requirements with MoSCoW priorities.}
\label{tab:func-req}
\end{table}

%==============================================================================
\section{Non-Functional Requirements}
%==============================================================================

\begin{table}[htbp]
\centering
\small
\setlength{\tabcolsep}{3pt}
\begin{tabularx}{\textwidth}{@{}ll>{\raggedright\arraybackslash}Xl@{}}
\toprule
\textbf{ID} & \textbf{Priority} & \textbf{Requirement} & \textbf{Verification} \\
\midrule
NFR1 & Must & The system must be reproducible: deterministic seeds, explicit artifact paths, and canonical CSV schemas. & Repeated runs; seed logging \\
\midrule
NFR2 & Must & The system must train on CPU without requiring GPU resources. & Successful training on CPU-only machine \\
\midrule
NFR3 & Must & The system must not commit or expose real patient data in version-controlled artifacts. & \texttt{.gitignore} rules; repository inspection \\
\midrule
NFR4 & Should & The system should handle datasets with up to 2 million rows without exceeding available memory. & Memory-efficient data iterator \\
\midrule
NFR5 & Should & The hyperparameter sweep should be resumable, skipping previously completed configurations. & Sweep rerun behaviour \\
\midrule
NFR6 & Should & The system should include a unit test suite with $>$50 tests covering core modules. & \texttt{pytest} output \\
\midrule
NFR7 & Could & The system could cache transformed data matrices to avoid redundant preprocessing across sweep runs. & Cache hit/miss logging \\
\bottomrule
\end{tabularx}
\caption{Non-functional requirements with MoSCoW priorities.}
\label{tab:nonfunc-req}
\end{table}

%==============================================================================
\section{Specification}
%==============================================================================

\subsection{Data Specification}

\begin{itemize}\tightlist
    \item \textbf{Input}: MIMIC-III hourly ICU tables produced by a MIMIC-Extract style pipeline, containing mixed-type features (continuous vital signs, binary intervention flags, categorical codes) and static demographics.
    \item \textbf{Autoregressive format}: Each row represents one patient-hour, with target columns (current-hour dynamic features) and condition columns (static covariates plus lagged dynamic features).
    \item \textbf{Hour-0 format}: One row per patient containing the joint static and initial dynamic state vector.
    \item \textbf{Output}: Synthetic CSV files in the same schema as the autoregressive input, decodable via the fitted preprocessor.
\end{itemize}

\subsection{Model Specification}

\begin{itemize}\tightlist
    \item \textbf{Backbone}: ForestFlow (all-continuous flow matching with XGBoost regressors) or HS3F (heterogeneous sequential generation with flow regressors for continuous and XGBoost classifiers for categorical features).
    \item \textbf{Time discretisation}: Configurable number of time levels $n_t$ (default: 20).
    \item \textbf{Noise multiplier}: Configurable $n_{\text{noise}}$ controlling training data augmentation (default: 50).
    \item \textbf{XGBoost parameters}: Number of estimators, max depth, and learning rate, all configurable via YAML.
    \item \textbf{Two-stage architecture}: IID hour-0 model plus conditional autoregressive model, trained independently but evaluated jointly in the sweep.
\end{itemize}

\subsection{Evaluation Specification}

\begin{itemize}\tightlist
    \item \textbf{Utility}: Multi-task TSTR with logistic regression baselines reporting accuracy, F1, and ROC-AUC per task.
    \item \textbf{Fidelity}: Per-feature KS statistics, correlation matrix Frobenius norm, clinical range violation percentage.
    \item \textbf{Privacy}: DCR median and match rate, baseline and overfitting protection scores, DOMIAS-style MIA ROC-AUC and attacker advantage.
    \item \textbf{Sweep}: Grid search over $(n_t, n_{\text{noise}})$ with results logged to a canonical CSV with schema enforcement.
\end{itemize}

%==============================================================================
\section{Known Limitations and Constraints}
%==============================================================================

\begin{itemize}\tightlist
    \item \textbf{Data access}: MIMIC-III requires credentialing and a signed data use agreement; the system cannot be tested without this access.
    \item \textbf{GitHub API rate limits}: Not applicable (no external API dependency), but sweep runtime scales as $O(n_t \times n_{\text{noise}} \times d_x)$ per configuration.
    \item \textbf{Baseline scope constraint}: A CTGAN/CRGAN baseline is included as a lower-bound/reference comparator, but it is not implemented as a fully autoregressive per-row conditioned generator. Extending it to a matched autoregressive architecture is a substantial design change and is outside this thesis scope.
    \item \textbf{Single dataset}: All requirements are specified and verified against MIMIC-III; portability to other clinical datasets is not guaranteed.
    \item \textbf{No real-time or interactive generation}: The system operates in batch mode via command-line scripts.
\end{itemize}
